\documentclass[a4paper,14pt]{extarticle}

% ── Кодировка и язык ──────────────────────────────────────────
\usepackage[T2A]{fontenc}
\usepackage[utf8]{inputenc}
\usepackage[russian]{babel}

% ── Поля ──────────────────────────────────────────────────────
\usepackage[left=30mm,right=15mm,top=20mm,bottom=20mm]{geometry}

% ── Оформление ────────────────────────────────────────────────
\usepackage{setspace}
\onehalfspacing
\usepackage{indentfirst}
\setlength{\parindent}{1.25cm}

% ── Математика ────────────────────────────────────────────────
\usepackage{amsmath,amssymb}

% ── Таблицы и графика ─────────────────────────────────────────
\usepackage{booktabs}
\usepackage{graphicx}
\usepackage{caption}
\usepackage{subcaption}

% ── Листинги ──────────────────────────────────────────────────
\usepackage{listings}

% ── Гиперссылки ───────────────────────────────────────────────
\usepackage[hidelinks]{hyperref}

% ── Библиография ──────────────────────────────────────────────
\usepackage[numbers,sort&compress]{natbib}
\bibliographystyle{unsrtnat}

% ── Секции без нумерации в оглавлении ─────────────────────────
\usepackage{titlesec}

% ══════════════════════════════════════════════════════════════
\begin{document}

% ── Титульный лист (заглушка) ─────────────────────────────────
\begin{titlepage}
\begin{center}
\vspace*{2cm}
{\large Министерство науки и высшего образования Российской Федерации}\\[0.5cm]
{\large Название университета}\\[0.5cm]
{\large Факультет / Институт}\\[0.5cm]
{\large Кафедра}\\[3cm]

{\Large\bfseries ВЫПУСКНАЯ КВАЛИФИКАЦИОННАЯ РАБОТА\\[0.3cm]
(бакалаврская работа)}\\[1.5cm]

{\large на тему:}\\[0.5cm]
{\Large\bfseries <<Исследование эффективности применения агентских систем\\
в задаче построения научных обзоров>>}\\[3cm]

\begin{flushright}
Выполнил: студент группы XXX\\
Фамилия~И.~О.\\[0.5cm]
Научный руководитель:\\
уч. степень, уч. звание\\
Фамилия~И.~О.
\end{flushright}

\vfill
{\large Город --- 2026}
\end{center}
\end{titlepage}

% ── Оглавление ────────────────────────────────────────────────
\tableofcontents
\newpage

% ══════════════════════════════════════════════════════════════
%  ВВЕДЕНИЕ
% ══════════════════════════════════════════════════════════════
\section*{Введение}
\addcontentsline{toc}{section}{Введение}

Объём научных публикаций растёт экспоненциально: по оценкам Bornmann и Mutz~\cite{bornmann2015growth}, среднегодовой прирост числа научных статей превышает 8\%.  Особенно стремительная динамика наблюдается в области искусственного интеллекта и обработки естественного языка, где только за 2023--2024 годы на платформе arXiv было размещено более 100~тысяч препринтов, связанных с большими языковыми моделями~\cite{wang2024autosurvey}. В таких условиях подготовка качественного обзора литературы~--- задача, традиционно требующая от исследователя недель систематической работы~--- становится узким местом научного процесса.

Появление больших языковых моделей (LLM) открыло возможность автоматизации этого процесса. За последние два года были предложены агентские системы, способные генерировать структурированные научные обзоры с цитированием источников: AutoSurvey~\cite{wang2024autosurvey}, SurveyForge~\cite{yan2025surveyforge}, STORM~\cite{shao2024storm}, SurveyGen-I~\cite{surveygen_i2025} и коммерческие решения, такие как Perplexity Deep Research. Эти системы применяют различные архитектурные подходы~--- от retrieval-augmented generation (RAG)~\cite{lewis2020rag} до мультиагентных pipeline с итеративным уточнением~--- и демонстрируют результаты, приближающиеся к человеческому уровню по ряду поверхностных характеристик.

Однако существующие метрики оценки качества генерируемых обзоров~--- ROUGE-L~\cite{lin2004rouge}, BERTScore~\cite{zhang2020bertscore}, G-Eval~\cite{liu2023geval}~--- были разработаны для коротких текстов (аннотаций, резюме) и не учитывают ключевые аспекты, отличающие хороший научный обзор от плохого. Они не измеряют внутреннюю согласованность длинного документа (наличие противоречий и повторов между секциями), не анализируют профиль фактологических ошибок по типам утверждений и не оценивают экспертные свойства письма~--- критичность, компаративность, эпистемическую модальность. Специализированный бенчмарк SurGE~\cite{su2025surge} предлагает четыре измерения оценки (полнота, точность цитирования, структура, качество содержания), но и он не покрывает перечисленные аспекты.

\textbf{Целью} настоящей работы является разработка и валидация набора автоматических метрик для оценки качества научных обзоров, сгенерированных LLM-based агентскими системами, и проведение сравнительного анализа существующих систем генерации на предложенных метриках.

Для достижения поставленной цели решаются следующие \textbf{задачи}:
\begin{enumerate}
    \item Провести анализ существующих систем генерации научных обзоров и методов оценки их качества.
    \item Разработать метрики межсекционной согласованности текста, включая автоматическое обнаружение противоречий и повторов между секциями обзора.
    \item Разработать метрики профилирования фактологической точности, позволяющие анализировать распределение ошибок цитирования по типам утверждений.
    \item Разработать метрики экспертных свойств научного письма: критичности, компаративности, формулирования открытых вопросов и эпистемической модальности.
    \item Валидировать предложенные метрики через экспертную разметку и оценку согласия автоматического классификатора с человеческими аннотаторами.
    \item Провести сравнительный анализ систем генерации обзоров (SurveyForge, SurveyGen-I, Perplexity Deep Research) на бенчмарке SurGE с использованием предложенных метрик.
\end{enumerate}

Работа состоит из четырёх глав. В главе~1 представлен обзор литературы по системам генерации обзоров, методам декомпозиции текста и существующим метрикам оценки. Глава~2 содержит формальную постановку задачи и описание исследовательских вопросов. В главе~3 подробно описаны предложенные метрики и архитектура вычислительного pipeline. Глава~4 посвящена экспериментальной оценке, валидации метрик и анализу результатов.

\newpage

% ══════════════════════════════════════════════════════════════
%  ГЛАВА 1. ОБЗОР ЛИТЕРАТУРЫ
% ══════════════════════════════════════════════════════════════
\section{Обзор литературы}

\subsection{Системы автоматической генерации научных обзоров}

Задача автоматической генерации научных обзоров привлекает растущее внимание исследовательского сообщества. Современные подходы к её решению можно разделить на три архитектурных парадигмы: pipeline-системы с однократной генерацией, системы с итеративным уточнением и коммерческие решения с закрытой архитектурой.

\textbf{Pipeline-системы.} AutoSurvey~\cite{wang2024autosurvey} стал одной из первых комплексных систем автоматической генерации обзоров. Система реализует четырёхэтапный pipeline: (1)~поиск релевантных публикаций в базе из 530~тысяч статей arXiv с помощью embedding-based retrieval; (2)~генерация структурированного плана обзора; (3)~параллельное написание подсекций специализированными LLM; (4)~интеграция и итеративное уточнение текста. Экспериментальная оценка на 20 темах из области LLM показала, что AutoSurvey достигает качества, сопоставимого с человеческими обзорами, при стоимости генерации около \$1.2 и времени менее 3~минут на обзор. Работа опубликована на NeurIPS~2024.

SurveyForge~\cite{yan2025surveyforge} развивает эту парадигму, вводя два ключевых улучшения. Во-первых, система анализирует логическую структуру существующих человеческих обзоров для эвристического построения плана, а не генерирует план <<с нуля>>. Во-вторых, SurveyForge использует агент навигации по научной литературе (scholar navigation agent) с механизмом памяти для более качественного извлечения релевантных публикаций. Авторы также предлагают бенчмарк SurveyBench, включающий 100 человеческих обзоров для сравнительной оценки. Работа принята на ACL~2025.

\textbf{Системы с итеративным уточнением.} STORM~\cite{shao2024storm}, разработанная в Стэнфордском университете, фокусируется на фазе pre-writing. Система моделирует исследовательский процесс через симуляцию диалогов между <<писателями>> с различными перспективами и <<экспертом по теме>>, опирающимся на интернет-источники. Собранная информация используется для построения структурированного плана, на основе которого генерируется статья в стиле Wikipedia. Хотя STORM ориентирована на энциклопедические статьи, а не на научные обзоры, её архитектурные решения (мультиперспективный поиск, кураторство информации) оказали значительное влияние на последующие работы.

SurveyGen-I~\cite{surveygen_i2025} представляет подход с эволюционирующим планом: в отличие от систем с фиксированным outline, план обзора обновляется по мере генерации содержания. Система построена на фреймворке LangGraph и использует memory-guided writing для поддержания согласованности текста. По результатам авторов, SurveyGen-I превосходит AutoSurvey и SurveyForge по всем пяти измерениям качества содержания.

\textbf{Коммерческие решения.} Системы Perplexity Deep Research и аналогичные продукты представляют отдельный класс решений с закрытой архитектурой. Они, как правило, комбинируют web-поиск в реальном времени с генерацией структурированного ответа, но не предоставляют доступа к внутреннему устройству pipeline. Включение таких систем в сравнительный анализ позволяет оценить разрыв между академическими и коммерческими решениями.


\subsection{Существующие метрики оценки качества генерируемых текстов}

Метрики оценки качества генерируемых текстов можно разделить на три категории: метрики лексического сходства, нейросетевые метрики и LLM-based оценки.

\textbf{Метрики лексического сходства.} ROUGE-L~\cite{lin2004rouge} измеряет пересечение $n$-грамм между сгенерированным текстом и эталоном. Несмотря на широкое распространение, ROUGE плохо коррелирует с человеческими оценками для длинных текстов и не учитывает семантическое сходство.

\textbf{Нейросетевые метрики.} BERTScore~\cite{zhang2020bertscore} вычисляет cosine similarity контекстуализированных представлений токенов между гипотезой и эталоном, что позволяет учитывать семантическую близость при различной лексике. AlignScore~\cite{zha2023alignscore} предлагает унифицированную функцию выравнивания на основе RoBERTa, обученную на задаче NLI, для проверки фактологической согласованности утверждения с контекстом.

\textbf{LLM-based оценки.} G-Eval~\cite{liu2023geval} использует GPT-4 для оценки текста по нескольким измерениям (coherence, consistency, fluency, relevance) с формированием цепочки рассуждений. Парадигма <<LLM-as-Judge>>~\cite{zheng2024llm_judge} получила широкое распространение благодаря высокой корреляции с человеческими оценками и масштабируемости.

\textbf{Специализированные метрики для обзоров.} Бенчмарк SurGE~\cite{su2025surge} предлагает оценку сгенерированных обзоров по четырём измерениям: (1)~полнота покрытия (comprehensiveness) на уровне документов, секций и предложений; (2)~точность цитирования; (3)~структурное качество (Structure Quality Score, SQS); (4)~качество содержания (Content Quality Score, CQS). SurveyBench~\cite{sun2025surveybench} дополнительно оценивает <<богатство>> обзора (наличие таблиц, диаграмм) и формулирует quiz-based оценку глубины покрытия.

Несмотря на значительный прогресс, ни одна из существующих метрик не анализирует внутреннюю согласованность длинного документа (противоречия и повторы между секциями), не профилирует фактологические ошибки по типам утверждений и не оценивает экспертные свойства научного письма.


\subsection{Декомпозиция текста на атомарные утверждения}

Стратегия <<decompose-then-verify>> стала стандартным подходом к проверке фактологичности длинных текстов~\cite{min2023factscore}. Она предполагает разбиение текста на атомарные утверждения (claims), каждое из которых может быть верифицировано независимо.

FActScore~\cite{min2023factscore} предложил простой подход: LLM разбивает предложение на <<atomic facts>>, которые затем проверяются против базы знаний (Wikipedia). Метод получил широкое распространение, но имеет два ограничения: он предполагает, что все утверждения верифицируемы, и не обеспечивает деконтекстуализацию (каждое утверждение должно быть понятно в изоляции от контекста).

VeriScore~\cite{song2024veriscore} решает первую проблему, различая верифицируемые и неверифицируемые утверждения. Метод использует один LLM-вызов на предложение, объединяющий классификацию, декомпозицию и деконтекстуализацию. Человеческая оценка подтвердила, что утверждения, извлечённые VeriScore, более осмысленны, чем у конкурирующих методов, на восьми различных задачах генерации длинных текстов. VeriScore выпущен под лицензией Apache~2.0.

Claimify~\cite{metropolitansky2025claimify} представляет наиболее развитый четырёхстадийный pipeline: (1)~selection~--- фильтрация предложений, не содержащих верифицируемого контента; (2)~disambiguation~--- выявление и разрешение референциальной и структурной неоднозначности; (3)~decomposition~--- разбиение на атомарные деконтекстуализированные утверждения. Уникальной особенностью Claimify является стадия disambiguation: система идентифицирует случаи, когда правильная интерпретация предложения не может быть установлена из контекста, и помечает такие предложения как <<Cannot be disambiguated>>. По результатам оценки на датасете BingCheck (396 ответов), Claimify достигает 99.0\% entailment (статистически неотличимо от VeriScore~--- 99.2\%) и значительно превосходит альтернативы по element-level coverage (macro~F1 83.7\% против 62.5\% у VeriScore).

SAFE~\cite{wei2024safe} и DnD~\cite{wanner2024dnd} представляют промежуточные решения. SAFE расширяет декомпозер FActScore дополнительными инструкциями и выполняет деконтекстуализацию отдельной стадией. DnD объединяет декомпозицию и деконтекстуализацию в одном промпте, но показывает наименьшую долю entailed утверждений (89.1\%), что свидетельствует о частом искажении исходного текста.

В настоящей работе используется VeriScore в качестве основного декомпозера, поскольку он обеспечивает оптимальное соотношение качества извлечения и вычислительной стоимости: один LLM-вызов на предложение против семи у Claimify при статистически неотличимом entailment.


\subsection{Парадигма LLM-as-Judge}

Использование больших языковых моделей в качестве автоматических оценщиков (LLM-as-Judge) стало одной из ключевых парадигм оценки качества генерируемых текстов~\cite{zheng2024llm_judge}. Подход основан на том, что инструктированные LLM способны выносить суждения о качестве текста, сопоставимые с экспертными оценками, при значительно меньших затратах времени и ресурсов.

Zheng~и~др.~\cite{zheng2024llm_judge} систематически исследовали этот подход на бенчмарке MT-Bench и платформе Chatbot Arena, показав, что GPT-4 достигает корреляции более 80\% с человеческими предпочтениями. Авторы также выявили типичные проблемы LLM-as-Judge: position bias (предпочтение первого ответа в парном сравнении), verbosity bias (предпочтение более длинных ответов) и self-enhancement bias (модель предпочитает собственные ответы).

В контексте оценки научных обзоров парадигма LLM-as-Judge используется в нескольких ролях: (1)~классификация утверждений по типам; (2)~проверка наличия экспертных свойств (критичность, компаративность); (3)~определение эпистемической модальности; (4)~верификация противоречий между высказываниями. Ключевое преимущество подхода~--- возможность формулировать сложные оценочные рубрики, недоступные для стандартных NLI-моделей. Ключевое ограничение~--- необходимость валидации через сравнение с человеческой разметкой, без которой результаты LLM-судьи не могут считаться достоверными.


\subsection{Выводы и выявленные пробелы}

Проведённый анализ позволяет выявить следующие пробелы в существующих подходах к оценке качества сгенерированных научных обзоров.

Во-первых, \emph{отсутствие метрик внутренней согласованности}. Ни одна из рассмотренных метрик не анализирует межсекционные противоречия~--- ситуации, когда в разных частях обзора делаются несовместимые утверждения об одном предмете. Между тем, именно этот тип ошибок характерен для LLM-генераторов, обрабатывающих каждую секцию с ограниченным контекстом предыдущих.

Во-вторых, \emph{отсутствие профилирования ошибок цитирования по типам утверждений}. Существующие метрики (FActScore, ALCE) считают общую долю поддержанных утверждений, не различая, какие типы утверждений чаще оказываются неподтверждёнными. Это не позволяет идентифицировать специфические слабости генераторов: например, систему, которая точно воспроизводит информацию из абстрактов, но галлюцинирует методологические детали.

В-третьих, \emph{отсутствие оценки экспертных свойств научного письма}. Качественный научный обзор содержит критический анализ, явные сравнения подходов, формулирование открытых вопросов и адекватную эпистемическую модальность (модуляцию уверенности в зависимости от степени обоснованности утверждения). Ни один из существующих фреймворков эти свойства не измеряет.

Настоящая работа направлена на восполнение указанных пробелов.


\newpage

% ══════════════════════════════════════════════════════════════
%  ГЛАВА 2. ПОСТАНОВКА ЗАДАЧИ
% ══════════════════════════════════════════════════════════════
\section{Постановка задачи}

\subsection{Формальная постановка}

Пусть $\mathcal{G} = \{G_1, G_2, \ldots, G_K\}$~--- множество систем генерации научных обзоров. Каждая система $G_k$ по заданной теме $t \in \mathcal{T}$ и набору ссылок $R_t = \{r_1, r_2, \ldots, r_n\}$ порождает обзор $S = G_k(t, R_t)$, представляющий собой структурированный текст, разбитый на секции $S = \{s_1, s_2, \ldots, s_m\}$.

Задача состоит в построении набора метрик $\mathcal{M} = \{M_1, M_2, \ldots, M_L\}$, каждая из которых отображает обзор $S$ в числовую оценку $M_i(S) \in [0, 1]$ (или $\mathbb{R}_+$ для некоторых метрик), характеризующую определённый аспект качества. Метрики должны удовлетворять следующим требованиям:

\begin{itemize}
    \item \textbf{Автоматизм:} вычисление $M_i(S)$ не требует участия эксперта (за исключением этапа валидации).
    \item \textbf{Дискриминативность:} метрика позволяет различать обзоры, написанные человеком, и сгенерированные различными системами.
    \item \textbf{Валидируемость:} результаты автоматического классификатора, лежащего в основе метрики, могут быть сопоставлены с экспертной разметкой через стандартные меры согласия (Cohen's~$\kappa$).
\end{itemize}


\subsection{Исследовательские вопросы}

Работа направлена на получение ответов на следующие исследовательские вопросы.

\textbf{RQ1:} Какие аспекты качества научных обзоров не покрываются существующими метриками, и можно ли их измерить автоматически с помощью LLM-as-Judge?

\textbf{RQ2:} Различаются ли LLM-генераторы обзоров по профилю фактологических ошибок~--- допускают ли они систематически больше ошибок в определённых типах утверждений (методологических, количественных, критических)?

\textbf{RQ3:} Воспроизводят ли LLM-генераторы эпистемическую осторожность, характерную для экспертного научного письма, или имитируют её поверхностно?


\subsection{Датасет SurGE}

В качестве основного бенчмарка используется SurGE (Survey Generation Evaluation)~\cite{su2025surge}~--- датасет для оценки генерации научных обзоров в области информатики. SurGE включает:

\begin{itemize}
    \item 205 тщательно отобранных эталонных обзоров с метаданными (заголовок, авторы, год публикации, аннотация, иерархическая структура, список ссылок);
    \item академический корпус из 1\,086\,992 статей с полями title, abstract, авторы, дата публикации и категория arXiv;
    \item автоматический фреймворк оценки по четырём измерениям: полнота покрытия, точность цитирования, структурное качество (SQS) и качество содержания (CQS).
\end{itemize}

Выбор SurGE обусловлен несколькими факторами. Во-первых, это крупнейший публично доступный бенчмарк для задачи генерации обзоров (205 эталонов против 20 в AutoSurvey и 100 в SurveyBench). Во-вторых, датасет включает полный корпус источников, что позволяет верифицировать утверждения обзора против абстрактов цитируемых статей. В-третьих, SurGE предоставляет собственный набор метрик, относительно которого можно позиционировать предложенные в данной работе метрики.


\newpage

% ══════════════════════════════════════════════════════════════
%  ГЛАВА 3. ПРЕДЛОЖЕННЫЕ МЕТРИКИ
% ══════════════════════════════════════════════════════════════
\section{Предложенные метрики}

В данной главе описывается набор метрик, разработанных для оценки трёх аспектов качества научного обзора, не покрываемых существующими фреймворками: межсекционной согласованности (группа~A), профиля фактологической точности (группа~B) и экспертных свойств научного письма (группа~C).


\subsection{Общая архитектура pipeline}

Вычисление метрик групп B и C требует предварительной декомпозиции текста обзора на атомарные утверждения. Для этого используется VeriScore~\cite{song2024veriscore}~--- метод, объединяющий классификацию предложения, декомпозицию на атомарные утверждения и деконтекстуализацию в одном LLM-вызове. Для каждого извлечённого утверждения сохраняется привязка к исходному предложению и параграфу обзора, что необходимо для корректной атрибуции источников в метрике~B.1.b.

Метрики группы~A работают непосредственно с предложениями обзора, без предварительной декомпозиции.


\subsection{Группа A: межсекционная согласованность}

Метрики данной группы оценивают внутреннюю согласованность текста обзора, анализируя пары предложений из различных секций.

\subsubsection{$M_\mathrm{contr}$: межсекционные противоречия}

Метрика обнаруживает пары предложений из разных секций обзора, которые делают несовместимые утверждения об одном и том же предмете. Вычисление осуществляется трёхстадийным pipeline.

\textbf{Стадия~1: предварительная фильтрация по семантической близости.} Для каждого предложения обзора вычисляется embedding с помощью SPECTER~\cite{cohan2020specter}~--- модели, обученной на задаче связывания научных статей через цитирования. Формируется матрица попарных cosine similarity, и отбираются межсекционные пары с similarity выше порога $\tau_\mathrm{sim}$:
\begin{equation}
    \mathrm{Cand}(S) = \{(p_i, p_j) : \mathrm{sec}(p_i) \ne \mathrm{sec}(p_j) \wedge \cos(\mathbf{e}_i, \mathbf{e}_j) > \tau_\mathrm{sim}\},
\end{equation}
где $\mathbf{e}_i$~--- SPECTER-embedding предложения $p_i$, $\tau_\mathrm{sim} = 0.6$~--- пороговое значение.

\textbf{Стадия~2: фильтрация по тематической связи (topic filter).} Для каждой пары из $\mathrm{Cand}(S)$ выполняется LLM-вызов с промптом: <<обсуждают ли два предложения один и тот же конкретный предмет (метод, результат, сущность)?>>. Пары, для которых ответ отрицательный, отбрасываются. Эта стадия заменяет NER-based предфильтрацию, которая требует корректной нормализации сущностей и плохо работает на научных текстах из-за доменного несоответствия существующих NER-моделей.

\textbf{Стадия~3: проверка противоречия (LLM-as-Judge).} Для оставшихся пар выполняется LLM-вызов с развёрнутой рубрикой, учитывающей специфику научного письма: переформулировки одного факта не считаются противоречиями, временная эволюция взглядов (<<ранние работы утверждали X, последующие показали Y>>) не является противоречием в голосе обзора, прагматические несовместимости (например, <<требует 8~A100 GPU в течение двух недель>> и <<вычислительно лёгкий подход>>) учитываются.

Итоговая метрика:
\begin{equation}
    M_\mathrm{contr}(S) = \frac{|\{(p_i, p_j) \in \mathrm{TopicOK}(S) : \mathrm{IsContr}(p_i, p_j) = 1\}|}{|\mathrm{TopicOK}(S)|},
\end{equation}
где $\mathrm{TopicOK}(S)$~--- множество пар, прошедших стадию~2.


\subsubsection{$M_\mathrm{rep}$: межсекционные повторы}

Метрика обнаруживает пары предложений из разных секций, передающих одну и ту же информацию. Pipeline двухстадийный:

\textbf{Стадия~1:} фильтрация по cosine similarity SPECTER-embeddings с порогом $\tau_\mathrm{rep} = 0.8$ (выше, чем для $M_\mathrm{contr}$, поскольку повторы по определению более лексически схожи).

\textbf{Стадия~2:} двунаправленная проверка через NLI-модель DeBERTa-v2-xlarge-mnli~\cite{he2021deberta}. Пара $(p_i, p_j)$ считается повтором тогда и только тогда, когда NLI-модель классифицирует и $(p_i, p_j)$, и $(p_j, p_i)$ как entailment. Двунаправленность необходима для отличия повтора от отношения уточнения (когда одно утверждение является частным случаем другого). Использование NLI вместо LLM-as-Judge обосновано тем, что повторы~--- это entailment на лексически близких парах, с чем NLI-модели справляются хорошо, в отличие от прагматических противоречий.

Формула:
\begin{equation}
    M_\mathrm{rep}(S) = \frac{|\{(p_i, p_j) : \mathrm{sim}(p_i, p_j) > \tau_\mathrm{rep} \wedge \mathrm{NLI}(p_i \to p_j) = \mathrm{ent} \wedge \mathrm{NLI}(p_j \to p_i) = \mathrm{ent}\}|}{|\{(p_i, p_j) : \mathrm{sim}(p_i, p_j) > \tau_\mathrm{rep}\}|}.
\end{equation}


\subsection{Группа B: профилирование фактологической точности}

\subsubsection{B.1.a: классификация утверждений}

Каждое атомарное утверждение, извлечённое декомпозером, классифицируется LLM-судьёй в одну из четырёх категорий:

\begin{itemize}
    \item \textbf{A~--- топические утверждения.} Утверждения, которые могут быть сделаны на основании только аннотации статьи: описание темы, проблемы или общего вклада работы.
    \item \textbf{B~--- методологические утверждения.} Утверждения, содержащие специфические детали о том, как работает метод, обычно из раздела Methods.
    \item \textbf{C~--- количественные утверждения.} Утверждения с конкретными числовыми значениями, как правило, из раздела Results.
    \item \textbf{D~--- критические и сравнительные утверждения.} Оценочные суждения, сравнения подходов, обсуждение ограничений~--- обычно из Discussion или Related Work.
\end{itemize}

При ошибке парсинга или API-вызова утверждение получает статус classification\_failed и исключается из знаменателя всех метрик, а не присваивается категория по умолчанию.


\subsubsection{B.1.b: условная корректность цитирования (CitCorrect$_k$)}

Для каждого утверждения определяются цитируемые источники через двухуровневую стратегию атрибуции:

\begin{enumerate}
    \item Парсинг inline-цитирований из исходного предложения, содержащего утверждение.
    \item Если предложение не содержит цитирований~--- fallback на парсинг цитирований из содержащего параграфа.
\end{enumerate}

Paragraph-level fallback необходим, поскольку в научном письме типичен паттерн, при котором ссылка указывается один раз в начале абзаца, а последующие предложения относятся к тому же источнику без повторной ссылки.

Для каждого атрибутированного утверждения проверяется его поддержка цитируемыми источниками с помощью AlignScore~\cite{zha2023alignscore} в режиме NLI (nli\_sp). Контекстом служит конкатенация абстрактов всех атрибутированных источников.

Условная корректность цитирования для категории $k$:
\begin{equation}
    \mathrm{CitCorrect}_k(S) = \frac{|\{a : \phi(a) = k \wedge \mathrm{Support}(a, e_a) \geq \theta\}|}{|\{a : \phi(a) = k \wedge \mathrm{status}(a) = \mathrm{ok}\}|},
\end{equation}
где $\phi(a) \in \{A, B, C, D\}$~--- категория утверждения, $e_a$~--- контекст из абстрактов атрибутированных источников, $\theta = 0.5$~--- порог AlignScore.

Утверждения без цитирований (no\_citations) и с ненайденными источниками (missing\_source) исключаются из знаменателя и репортируются отдельно.


\subsection{Группа C: экспертные свойства научного письма}

Метрики данной группы оценивают четыре свойства, характерные для качественного экспертного научного письма. Каждое свойство определяется независимым бинарным (или ординальным) классификатором, применяемым к атомарным утверждениям.


\subsubsection{$M_\mathrm{crit}$: критичность}

Доля утверждений, содержащих критический анализ: указание на ограничения, отрицательные результаты, противоречия, компромиссы (trade-offs). Гипотеза: качественный обзор содержит больше критических суждений, чем поверхностный пересказ абстрактов.

\begin{equation}
    M_\mathrm{crit}(S) = \frac{|\{a : \mathrm{IsCritical}(a) = 1\}|}{|\{a\}|}.
\end{equation}


\subsubsection{$M_\mathrm{comp}$: компаративность}

Метрика расщеплена на две компоненты. $M_\mathrm{comp\text{-}total}$~--- доля утверждений, содержащих явное сравнение нескольких работ или подходов:
\begin{equation}
    M_\mathrm{comp\text{-}total}(S) = \frac{|\{a : \mathrm{IsComparative}(a) = 1\}|}{|\{a\}|}.
\end{equation}

$M_\mathrm{comp\text{-}valid}$~--- доля валидных сравнений среди компаративных утверждений:
\begin{equation}
    M_\mathrm{comp\text{-}valid}(S) = \frac{|\{a : \mathrm{IsComparative}(a) = 1 \wedge \mathrm{IsValid}(a) = 1\}|}{|\{a : \mathrm{IsComparative}(a) = 1\}|}.
\end{equation}

Сравнение считается невалидным, если сравниваемые методы оценивались на разных датасетах, с использованием различных метрик или при несопоставимых экспериментальных условиях.


\subsubsection{$M_\mathrm{open}$: формулирование открытых вопросов}

Доля утверждений, формулирующих нерешённые исследовательские вопросы, направления будущей работы или области, требующие дальнейшего изучения:
\begin{equation}
    M_\mathrm{open}(S) = \frac{|\{a : \mathrm{IsOpenQuestion}(a) = 1\}|}{|\{a\}|}.
\end{equation}


\subsubsection{$M_\mathrm{mod}$: эпистемическая модальность}

Метрика оценивает распределение утверждений по пяти уровням эпистемической уверенности:

\begin{enumerate}
    \item \textbf{Категоричные}~--- без хеджирования: <<is>>, <<does>>, <<demonstrates>>.
    \item \textbf{Сильные}~--- мягкие квалификаторы: <<generally>>, <<typically>>, <<has been shown to>>.
    \item \textbf{Умеренные}~--- явное хеджирование: <<often>>, <<tends to>>, <<can>>.
    \item \textbf{Слабые}~--- сильное хеджирование: <<may>>, <<might>>, <<could>>, <<suggests>>.
    \item \textbf{Эксплицитная неопределённость}~--- <<remains unclear>>, <<is debated>>, <<open question>>.
\end{enumerate}

Итоговая метрика~--- энтропия Шеннона распределения по уровням:
\begin{equation}
    M_\mathrm{mod}(S) = -\sum_{l=1}^{5} p_l \log_2 p_l,
\end{equation}
где $p_l$~--- доля утверждений на уровне $l$. Более высокая энтропия указывает на большее разнообразие модальностей, что характерно для экспертного письма, модулирующего степень уверенности в зависимости от обоснованности утверждения.


\subsection{Выбор моделей и параметров pipeline}

Для LLM-as-Judge во всех метриках используются открытые модели, доступные через OpenRouter API. Выбор модели определяется балансом качества и стоимости: для задач, требующих глубокого reasoning (проверка противоречий, оценка валидности сравнений), применяется модель с более высокими показателями на бенчмарках GPQA и MMLU-Pro; для простых классификационных задач~--- модель с лучшим соотношением скорости и цены.

Все промпты построены с учётом KV-cache оптимизации: статическая часть (рубрика, формат ответа) размещается в начале, переменная часть (текст утверждений) --- в конце. Это позволяет провайдеру переиспользовать кэш для статического префикса при обработке множества утверждений.

Параллелизация LLM-вызовов осуществляется через ThreadPoolExecutor с конфигурируемым числом параллельных запросов. Кэширование результатов на уровне отдельных вызовов (diskcache по хэшу от входных данных) обеспечивает устойчивость к обрывам и возможность возобновления прогона без повторных API-вызовов.


\newpage

% ── Список литературы ─────────────────────────────────────────
\bibliography{references}

\end{document}